\section{Motivation}
The South African mining industry continues to experience a disproportionate burden of occupational lung disease, with miners routinely exposed to silica dust, diesel particulates, and bio-aerosols that drive high incidence rates of silicosis, tuberculosis (TB), and pneumonia. Radiographic screening remains the primary triage tool for these co-morbidities, yet many clinics operate with ageing X-ray systems that produce heterogeneous image characteristics and low signal-to-noise ratios. These conditions fragment the available labelled datasets, leaving diagnostic models under-trained on the very edge cases---mixed silico-tuberculosis presentations---that matter most for clinical triage in mining communities. A data-efficient method that can faithfully expand these underrepresented classes is therefore critical for producing classifiers that generalise beyond the limited samples gathered in specific hospitals or machine settings.

Diffusion-based generative models provide a principled avenue for synthesising realistic chest radiographs without the heavy annotation burden of supervised acquisition. More specifically, recent advances in compositional diffusion models enable controllable synthesis at inference time by factorising complex scenes into semantically meaningful components. Applying this paradigm to chest X-rays creates an opportunity to model silicosis-associated reticulations and TB-driven cavitations as separable but interacting latent concepts. By conditioning on learned pathology factors and recombining them in novel configurations, we can synthesise clinically plausible co-morbidity cases that are rarely captured in real-world datasets. Crucially, inference-time compositionality avoids the need to retrain the diffusion backbone whenever a new pathology combination must be explored, dramatically lowering the barrier to iterative experimentation with occupational lung disease phenotypes.

Beyond data augmentation, controllable synthesis offers a mechanism for hypothesis-driven interrogation of latent disease factors. Radiologists and epidemiologists can request targeted perturbations---such as varying silica exposure severity while holding TB lesion distribution constant---to probe diagnostic cues and improve human understanding of disease progression. Such capabilities align with South Africa's public health priorities by enabling rapid scenario analysis, supporting tele-radiology workflows in remote mining towns, and ultimately improving early intervention outcomes.